\section{Distribución normal}
En estadística, cuando ordenamos los resultados de un fenómeno educativo y la mayoría de los
valores se concentra alrededor de un centro, disminuyendo de forma pareja hacia ambos extremos,
aparece la forma de montaña conocida como distribución normal o campana de Gauss. Es el modelo
matemático más importante para analizar fenómenos educativos y sociales, ya que se define por
completo con solo dos parámetros. En esta sección aplicaremos este modelo sobre la variable
\Variable{english\_score}, cuyo análisis reveló un comportamiento notablemente parejo.

\subsection{Parámetros de la distribución}
Al ordenar las notas de inglés, notamos que la mayoría de los estudiantes se concentra alrededor
de una nota central y que, a medida que nos alejamos de ese centro (hacia notas excelentes o notas
muy bajas), la cantidad de alumnos disminuye de forma pareja. La campana se define completamente
con los siguientes dos datos:
\begin{itemize}
	\item \textbf{La media ($\mu = 82{,}54$):} El centro exacto de la campana, que marca el
	rendimiento promedio del grupo en inglés.
	\item \textbf{La desviación estándar ($\sigma = 3{,}85$):} La medida que nos dice qué tan
	dispersas o qué tan agrupadas están las notas respecto al promedio.
\end{itemize}

\PalabraClave{Conclusiones clave para la gestión educativa:}
\begin{itemize}
	\item \textbf{Rendimiento general sólido y homogéneo:} El promedio del grupo se sitúa en
	82,54 puntos con una dispersión baja de 3,85 puntos. Esto demuestra que el nivel de inglés de
	los alumnos es, en general, alto y bastante parejo, sin registrarse desniveles extremos ni
	alarmantes en la mayoría del aula.
	\item \textbf{Identificación temprana de alumnos en riesgo:} Gracias al modelo, determinamos
	que existe un 25,46\% de probabilidad de que un estudiante obtenga una nota menor a 80
	puntos. Esta métrica permite a la institución anticipar qué porcentaje del grupo podría
	requerir tutorías o apoyo pedagógico.
	\item \textbf{Criterio de becas justo y transparente:} Para cumplir con la condición de
	premiar únicamente al 23\% de los mejores promedios, la matemática definió que la nota de
	corte exacta debe ser 85,50 puntos. Utilizar este umbral garantiza un criterio de selección
	objetivo, basado estrictamente en la distribución del mérito del grupo.
	\item \textbf{Concentración del rendimiento:} Casi la mitad de los estudiantes (47,27\%) se
	concentra en el rango de calificaciones entre 77 y 83 puntos. Esto confirma la alta densidad
	de alumnos con un desempeño estándar-bueno, característico de una distribución normal bien
	balanceada.
\end{itemize}

\AnalisisDescriptivo
\begin{TablaDatosFlexible}{L{0.28\linewidth} C{0.16\linewidth} L{0.46\linewidth}}{\textbf{Medida} & \textbf{Valor} & \textbf{Interpretación}}

\TablaFilaFlexible{Cantidad de datos ($n$) & 111 & La muestra está compuesta por 111 estudiantes.}
\TablaFilaFlexible{Media aritmética ($\mu$) & 82,54 & Es el promedio central de la campana: en torno a este valor se concentra el rendimiento en inglés del grupo.}
\TablaFilaFlexible{Desviación estándar poblacional ($\sigma$) & 3,85 & Mide qué tan dispersas están las notas respecto al promedio, considerando los 111 alumnos como población.}
\TablaFilaFlexible{Desviación estándar muestral ($s$) & 3,87 & Resultado equivalente si se considera que los datos provienen de una muestra.}
\TablaFilaFlexible{Calificación mínima & 77 & La nota más baja del grupo se ubica bastante cerca del centro.}
\TablaFilaFlexible{Calificación máxima & 92 & La nota más alta dista menos de 10 puntos del promedio.}
\TablaFilaFlexible{Rango & 15 & La brecha total entre la nota máxima y la mínima es de apenas 15 puntos.}

\end{TablaDatosFlexible}

% Definición de la función de densidad normal para los gráficos de esta sección.
\pgfmathdeclarefunction{gauss}{3}{%
  \pgfmathparse{1/(#3*sqrt(2*pi))*exp(-((#1-#2)^2)/(2*(#3)^2))}%
}

\subsection{Preguntas}

\subsubsection{Pregunta 1: Probabilidad de nota menor a 80}

¿Cuál es la probabilidad de que un estudiante elegido al azar obtenga una calificación menor a
80 puntos?

\PalabraClave{Cómo se toman los datos:}
Primero estandarizamos la nota de corte a la escala estándar normal ((Z)), restando la media y
dividiendo por la desviación estándar poblacional:
\begin{equation*}
  Z = \frac{80 - \mu}{\sigma} = \frac{80 - 82{,}54}{3{,}85} = \frac{-2{,}54}{3{,}85} \approx -0{,}66
\end{equation*}

\paragraph{Respuesta:}
\[ P(X < 80) = P(Z < -0{,}66) = 0{,}2546 = 25{,}46\% \]

\paragraph{Interpretación:} La probabilidad de que un estudiante obtenga una calificación menor a
80 puntos es de 25,46\%. Es decir, hay aproximadamente una posibilidad en cuatro de que un alumno
elegido al azar quede por debajo de esa nota, lo que ayuda a anticipar la demanda de tutorías.

\TituloGrafico{Área bajo la curva para \texorpdfstring{$P(X < 80)$}{P(X menor a 80)}}
\begin{figure}[H]
    \centering
    \begin{tikzpicture}
        \begin{axis}[
            width=0.7\textwidth,
            height=6.5cm,
            xlabel={Nota de inglés (puntos)},
            ylabel={Densidad de probabilidad},
            xmin=70, xmax=95,
            ymin=0, ymax=0.12,
            xtick={70,75,80,85,90,95},
            ytick=\empty,
            axis lines=left,
            samples=200,
            clip=false,
        ]
            % Curva normal completa.
            \addplot[thick, blue!60!black, domain=70:95] {gauss(x,82.54,3.85)};
            % Región sombreada: P(X < 80).
            \addplot[draw=none, fill=blue!30, domain=70:80] {gauss(x,82.54,3.85)} \closedcycle;
            % Líneas de referencia.
            \draw[dashed, gray] (axis cs:82.54,0) -- (axis cs:82.54,0.11)
                node[above, black] {$\mu = 82{,}54$};
            \draw[dashed, blue!60!black] (axis cs:80,0) -- (axis cs:80,{gauss(80,82.54,3.85)});
            \node[above] at (axis cs:75,0.03) {$P = 25{,}46\%$};
        \end{axis}
    \end{tikzpicture}
    \caption{Distribución normal (\(\mu = 82{,}54\), \(\sigma = 3{,}85\)) con el área sombreada
    representando la probabilidad de obtener una nota menor a 80 puntos.}
    \label{fig:normal_menor_80}
\end{figure}

% --------------------------------------------------------------------------------------------------

\subsubsection{Pregunta 2: Probabilidad de nota mayor a 90}

¿Cuál es la probabilidad de que un estudiante elegido al azar obtenga una calificación mayor a
90 puntos?

\PalabraClave{Cómo se toman los datos:}
Nuevamente estandarizamos la nota de corte:
\begin{equation*}
  Z = \frac{90 - \mu}{\sigma} = \frac{90 - 82{,}54}{3{,}85} = \frac{7{,}46}{3{,}85} \approx 1{,}94
\end{equation*}

\paragraph{Respuesta:}
\[ P(X > 90) = 1 - P(Z < 1{,}94) = 1 - 0{,}9738 = 0{,}0262 = 2{,}62\% \]

\paragraph{Interpretación:} La probabilidad de que un estudiante obtenga una calificación mayor a
90 puntos es de apenas 2,62\%. Las notas sobresalientes son, por lo tanto, un fenómeno extremo
dentro de esta distribución: apenas dos o tres alumnos de cada cien alcanzan ese nivel.

\TituloGrafico{Área bajo la curva para \texorpdfstring{$P(X > 90)$}{P(X mayor a 90)}}
\begin{figure}[H]
    \centering
    \begin{tikzpicture}
        \begin{axis}[
            width=0.7\textwidth,
            height=6.5cm,
            xlabel={Nota de inglés (puntos)},
            ylabel={Densidad de probabilidad},
            xmin=70, xmax=95,
            ymin=0, ymax=0.12,
            xtick={70,75,80,85,90,95},
            ytick=\empty,
            axis lines=left,
            samples=200,
            clip=false,
        ]
            % Curva normal completa.
            \addplot[thick, blue!60!black, domain=70:95] {gauss(x,82.54,3.85)};
            % Región sombreada: P(X > 90).
            \addplot[draw=none, fill=blue!30, domain=90:95] {gauss(x,82.54,3.85)} \closedcycle;
            % Líneas de referencia.
            \draw[dashed, gray] (axis cs:82.54,0) -- (axis cs:82.54,0.11)
                node[above, black] {$\mu = 82{,}54$};
            \draw[dashed, blue!60!black] (axis cs:90,0) -- (axis cs:90,{gauss(90,82.54,3.85)});
            \node[above] at (axis cs:91.5,0.02) {$P = 2{,}62\%$};
        \end{axis}
    \end{tikzpicture}
    \caption{Distribución normal (\(\mu = 82{,}54\), \(\sigma = 3{,}85\)) con el área sombreada
    representando la probabilidad de obtener una nota mayor a 90 puntos.}
    \label{fig:normal_mayor_90}
\end{figure}

% --------------------------------------------------------------------------------------------------

\subsubsection{Pregunta 3: Probabilidad de nota entre 77 y 83}

¿Cuál es la probabilidad de que un estudiante elegido al azar obtenga una calificación entre 77
y 83 puntos?

\PalabraClave{Cómo se toman los datos:}
Estándarizamos ambos límites del intervalo:
\begin{equation*}
  Z_1 = \frac{77 - 82{,}54}{3{,}85} = \frac{-5{,}54}{3{,}85} \approx -1{,}44
  \qquad
  Z_2 = \frac{83 - 82{,}54}{3{,}85} = \frac{0{,}46}{3{,}85} \approx 0{,}12
\end{equation*}

\paragraph{Respuesta:}
\begin{equation*}
  \begin{split}
    P(77 < X < 83) &= P(-1{,}44 < Z < 0{,}12) = P(Z < 0{,}12) - P(Z < -1{,}44) \\
                    &= 0{,}5477 - 0{,}0750 = 0{,}4727 = 47{,}27\%
  \end{split}
\end{equation*}

\paragraph{Interpretación:} Elegiendo la nota de inglés de un estudiante al azar, la probabilidad
de que la misma esté entre 77 y 83 puntos es de 47,27\%. Casi la mitad del grupo se concentra en
ese estrecho rango de seis puntos, lo que confirma la alta densidad de alumnos con un desempeño
estándar-bueno.

\TituloGrafico{Área bajo la curva para \texorpdfstring{$P(77 < X < 83)$}{P(77 menor que X menor que 83)}}
\begin{figure}[H]
    \centering
    \begin{tikzpicture}
        \begin{axis}[
            width=0.7\textwidth,
            height=6.5cm,
            xlabel={Nota de inglés (puntos)},
            ylabel={Densidad de probabilidad},
            xmin=70, xmax=95,
            ymin=0, ymax=0.12,
            xtick={70,75,80,85,90,95},
            ytick=\empty,
            axis lines=left,
            samples=200,
            clip=false,
        ]
            % Curva normal completa.
            \addplot[thick, blue!60!black, domain=70:95] {gauss(x,82.54,3.85)};
            % Región sombreada: P(77 < X < 83).
            \addplot[draw=none, fill=blue!30, domain=77:83] {gauss(x,82.54,3.85)} \closedcycle;
            % Líneas de referencia.
            \draw[dashed, gray] (axis cs:82.54,0) -- (axis cs:82.54,0.11)
                node[above, black] {$\mu = 82{,}54$};
            \draw[dashed, blue!60!black] (axis cs:77,0) -- (axis cs:77,{gauss(77,82.54,3.85)});
            \draw[dashed, blue!60!black] (axis cs:83,0) -- (axis cs:83,{gauss(83,82.54,3.85)});
            \node[above] at (axis cs:80,0.03) {$P = 47{,}27\%$};
        \end{axis}
    \end{tikzpicture}
    \caption{Distribución normal (\(\mu = 82{,}54\), \(\sigma = 3{,}85\)) con el área sombreada
    representando la probabilidad de obtener una nota entre 77 y 83 puntos.}
    \label{fig:normal_entre_77_83}
\end{figure}

\Interpretacion{Los tres casos analizados muestran cómo la distribución normal traduce la nota de corte en una probabilidad precisa. El bajo desvío estándar (3,85 puntos) concentra a casi la mitad del grupo en el rango de 77 a 83 puntos y vuelve excepcional el rendimiento sobresaliente (mayor a 90), mientras que identifica a un cuarto del grupo como posible destinatario de apoyo pedagógico. El modelo confirma que el nivel de inglés del aula es alto y homogéneo, y ofrece a la institución umbrales objetivos tanto para las becas (85,50 puntos) como para las tutorías (80 puntos).}


